\documentclass[11pt,a4paper]{article}
\usepackage[utf8]{inputenc}
\usepackage[T5]{fontenc}
\usepackage[vietnamese]{babel}
\usepackage{amsmath,amssymb,amsthm}
\usepackage{geometry,booktabs,tabularx,longtable,array}
\usepackage{tikz}
\usetikzlibrary{arrows.meta,positioning}
\usepackage{xcolor,hyperref,xurl,fancyhdr}
\geometry{margin=2.3cm,headheight=15pt}
\definecolor{ink}{HTML}{17354A}
\definecolor{accent}{HTML}{177A80}
\hypersetup{colorlinks=true,linkcolor=ink,citecolor=accent,urlcolor=accent,
 pdftitle={ECDR: Báo cáo nghiên cứu và hiện trạng},pdfauthor={}}
\urlstyle{same}
\setlength{\parindent}{0pt}
\setlength{\parskip}{5pt}
\setlength{\emergencystretch}{3em}
\setlength{\tabcolsep}{5pt}
\renewcommand{\arraystretch}{1.18}
\newcolumntype{Y}{>{\raggedright\arraybackslash}X}
\newcommand{\R}{\mathbb{R}}
\newcommand{\Sset}{\mathcal{S}}
\newcommand{\Zset}{\mathcal{Z}}
\newcommand{\code}[1]{\texttt{\detokenize{#1}}}
\newcommand{\norm}[1]{\left\lVert #1\right\rVert}
\newtheorem{proposition}{Mệnh đề}
\pagestyle{fancy}
\fancyhf{}
\fancyhead[L]{\small\textcolor{ink}{ECDR | Báo cáo nghiên cứu}}
\fancyhead[R]{\small 22/10/2026}
\fancyfoot[C]{\thepage}
\title{\textcolor{ink}{Tái sử dụng quyết định có chứng nhận\\
theo sự kiện trong suy luận mạng neural\\trên chuỗi cảm biến}\\[8pt]
\large Event-Conditioned Certified Decision Reuse (ECDR)\\
\normalsize Báo cáo hiện trạng và đề cương phương pháp}
\author{}
\date{Ngày cập nhật: 22 tháng 10 năm 2026}

\begin{document}
\maketitle
\begin{abstract}
ECDR nghiên cứu khả năng tránh chạy lại một neural classifier đã cố định trên
mỗi cửa sổ cảm biến bằng cách tái sử dụng quyết định tại một trusted anchor.
Việc chấp nhận reuse đòi hỏi đầu vào mới thuộc đúng miền được chứng nhận và
mọi class-pair logit margin có lower bound hợp lệ, dương trên toàn miền.
Báo cáo xác định bài toán, câu hỏi nghiên cứu, domain contract, điều kiện đủ
của certificate, kiến trúc runtime dự kiến và phương pháp đánh giá.
Nguồn thực nghiệm hiện có là các artifact HAPT đã lưu: 17.453 cửa sổ,
backbone sáu lớp có 282.918 tham số, test macro-F1 của checkpoint chọn bằng
validation là 0,9354 trên 2.874 pure windows. Mean forward CPU được báo cáo
là 0,2511 ms; đây chưa phải latency end-to-end. Các số liệu này chưa được
tái lập bằng chạy model trong lượt biên soạn. Artifact ghi nhận các run ngày
15/10/2026; ngày cập nhật báo cáo không phải ngày thực nghiệm mới. Verifier, certified reuse gate
và net saving chưa có kết quả. Đóng góp ECDR vì vậy còn là giả thuyết cần
kiểm nghiệm, với G3--G6 đang chờ thực hiện.
\end{abstract}
\textbf{Từ khóa:} certified decision reuse; time-series classification; neural
network verification; CROWN; perturbation region; trusted anchor; runtime cost.

\tableofcontents
\newpage

\section{Bài toán và động lực nghiên cứu}
Một hệ thống phân loại chuỗi cảm biến thường dựng cửa sổ rồi chạy đầy đủ
backbone cho từng cửa sổ. Các đầu vào liên tiếp có thể chứa redundancy, nhưng
chồng lấn hoặc gần nhau theo một khoảng cách không đủ để bảo đảm classifier
giữ quyết định. Khi chỉ dùng heuristic similarity để reuse, hệ thống có thể
trả một nhãn khác với nhãn mà backbone sẽ tính cho query hiện tại.

ECDR đặt mục tiêu kiểm tra điều kiện đủ trước khi reuse. Đối tượng được bảo
toàn là \emph{quyết định của reference pipeline}; nhãn hoạt động thực tế vẫn
có thể bị backbone phân loại sai. Chứng nhận cùng quyết định không tạo ra
guarantee correctness với ground truth, không tái tạo toàn bộ logits và
không mở rộng context window hoặc bộ nhớ của LLM.

Giá trị ứng dụng cần được đo thay vì giả định. Backbone hiện tại có chi phí
forward nhỏ trên CPU đã benchmark. Chưa có deployment, kết quả trên điện
thoại, đo energy hay bằng chứng rằng một stream đang vi phạm deadline.

\section{Câu hỏi nghiên cứu và giả thuyết}
\textbf{RQ1 -- câu hỏi chính.}
Miền biến thiên được chọn theo event và class-pair bounds có thể chứng nhận
giữ nguyên quyết định backbone thường xuyên hơn global Lipschitz bounds,
trong khi tổng chi phí đủ thấp để có net inference saving hay không?

\begin{itemize}
\item \textbf{RQ2 -- soundness:} cần domain, membership, numerical assumptions
và identity nào để mọi query được chấp nhận đều giữ quyết định?
\item \textbf{RQ3 -- effectiveness:} event-conditioned structure có cải thiện
certified coverage/reuse khi so sánh ở budget hoặc coverage phù hợp?
\item \textbf{RQ4 -- systems cost:} mức reuse và certificate amortization nào
bù được event processing, membership, fallback và rebuild?
\end{itemize}

\textbf{H1 (điều kiện toán học):} membership đúng và lower margins hợp lệ,
dương trên toàn miền là điều kiện đủ cho decision preservation.
\textbf{H2 (thực nghiệm, mở):} cấu trúc miền cải thiện certified reuse ở ít
nhất một operating regime có so sánh công bằng.
\textbf{H3 (thực nghiệm, mở):} có operating regime mà tổng chi phí thấp hơn
full inference. H2/H3 chưa có kết quả; không đặt target cứng giảm 50\% latency.

\section{Phạm vi, dữ liệu và các đối tượng}
\subsection{Window, event và region}
Một time-series window là một điểm đầu vào
\(x_t\in\R^{6\times128}\), tương đương vector trong \(\R^{768}\).
Mỗi window chứa 128 thời điểm của sáu kênh accelerometer/gyroscope.
Ở 50 Hz, độ dài quy ước là 2,56 giây; stride 64 tương ứng 1,28 giây.
Các thông số này phù hợp cách dựng cửa sổ của project và mô tả HAPT
\cite{hapt}. Project sử dụng raw recordings để dựng lại chronology; số
17.453 cửa sổ của project không phải số 10.929 feature records trên trang UCI.

\begin{figure}[ht]
\centering
\begin{tikzpicture}[x=0.065cm,y=0.65cm,font=\small]
\draw[-{Latex}] (0,0)--(194,0) node[right]{mẫu};
\fill[accent!18] (0,0.4) rectangle (128,1.0);
\draw[accent,thick] (0,0.4) rectangle (128,1.0);
\fill[ink!14] (64,1.5) rectangle (192,2.1);
\draw[ink,thick] (64,1.5) rectangle (192,2.1);
\node at (64,0.7){\(x_t:\ 128\) thời điểm/kênh};
\node at (128,1.8){\(x_{t+1}:\ 128\) thời điểm/kênh};
\draw[dashed] (64,0)--(64,2.3);
\draw[dashed] (128,0)--(128,2.3);
\foreach \x/\t in {0/0,64/64,128/128,192/192}
 {\draw (\x,0.07)--(\x,-0.07) node[below]{\t};}
\end{tikzpicture}
\caption{Minh họa hình học của sliding windows; không phải đồ thị tín hiệu
đo được. Hai cửa sổ chia sẻ 64 thời điểm, nhưng các mẫu đó dịch vị trí trong
tensor CNN. Overlap không tự suy ra decision agreement.}
\end{figure}

Event descriptor \(e\) mô tả biến thiên và giúp selector chọn template/budget.
Selector online hiện chưa được triển khai. Nhãn ground-truth dùng trong
audit không có nghĩa runtime biết đúng hoạt động.
Candidate region \(\Sset_e(x^r)\) là một \emph{tập} đầu vào gắn với anchor,
không phải một đoạn thời gian. Chỉ khi có certificate hợp lệ trên toàn miền
mới gọi đó là certified region.

\subsection{Phạm vi classification và model}
Backbone cố định hiện tại có \(K=6\) basic activities. K=12 pure-only đang HOLD
vì transition support không đạt tiêu chí project.
Ba khối Conv1d có channels \(6\to32\to64\to64\), kernels \(5/3/3\);
mỗi conv có BatchNorm và ReLU, MaxPool sau hai conv đầu. Head là
Flatten 2048--Linear 128--ReLU--Linear 6. Tổng tham số là 282.918.
Đây là kiến trúc thực trong source, thay cho proposal no-BN khoảng 24k cũ.
Khả năng tạo bounds cho graph thực tế còn cần kiểm trong G3.

\section{Công trình liên quan và khoảng trống}
LipCache nghiên cứu certified caching cho image service; bản PDF đã đọc
giới hạn certificate ở phía GuardNet, không bảo đảm MainNet trên toàn reuse
ball \cite{lipcache}. Wei và Liu tái sử dụng computation cho online
verification dưới domain/network updates \cite{wei}; đây là mục tiêu
verification, không tự tương đương việc tránh forward inference.
NoScope và DeltaCNN khai thác temporal structure cho conditional hoặc sparse
inference \cite{noscope,deltacnn}. Các nhánh đó đặt ra prior art cho temporal
reuse, nhưng có guarantee và phạm vi khác.

CROWN và automatic LiRPA cung cấp nền tảng tính bounds
\cite{crown,lirpa}. API auto\_LiRPA hỗ trợ specification matrix \(C\) ánh xạ
logits thành classification margins \cite{api}; tính năng API chưa chứng
minh integration hiện tại đúng hoặc đủ nhanh. STBP và time-series smoothing
là prior art cho structured perturbations/certification
\cite{stbp,franco}.

\textbf{Candidate contribution của ECDR} là tổ hợp cụ thể: direct-backbone
decision preservation, miền được chọn theo event, sound membership,
trusted-anchor lifecycle và end-to-end cost accounting. Chưa chứng minh
tổ hợp này mới. Không nhận first temporal reuse, first structured verifier,
first class-pair margin hay net-saving như kết quả đã đạt.

Coverage của literature review còn giới hạn: có reading list P01--P23,
claim ledger và pilot note P04; chưa sàng lọc toàn corpus theo systematic
search. Metadata check hoặc đọc một đoạn không tương đương full-paper review.

\section{Domain contract và formulation}
\subsection{Raw và normalized coordinates}
Gọi \(x\in\R^d\), \(d=768\), là raw window đã flatten.
Preprocessing train-fit cố định là
\begin{equation}
 z=N(x)=A(x-m),\qquad
 A=\operatorname{diag}(\sigma_1^{-1},\ldots,\sigma_d^{-1}),\qquad
 F(x)=f_\theta(N(x)).
\end{equation}
Mean/std theo channel được lặp theo time coordinate; std trong source có
floor \(10^{-6}\). \(F\) là raw-input reference pipeline, còn \(f_\theta\)
nhận normalized input. Model weights, BatchNorm inference state và
normalization phải được pin cùng certificate.

Ở G3 đầu tiên, định nghĩa normalized diagonal candidate
\begin{equation}\label{eq:diag}
 \Zset_e(z^r)=
 \{z^r+B_eu:\norm{u}_\infty\le1\},\quad
 B_e=\operatorname{diag}(b_1,\ldots,b_d),\quad b_i\ge0,
\end{equation}
với \(z^r=N(x^r)\). Raw preimage tương ứng là
\(\Sset_e(x^r)=N^{-1}(\Zset_e(z^r))\).
Do đó không áp raw budget trực tiếp lên normalized model.
Đây là contract cho prototype dự kiến; budget concrete và integration chưa
được triển khai hoặc freeze thành kết quả thực nghiệm.

\subsection{Certificate trên toàn miền}
Tại trusted anchor đã chạy full model, đặt
\(a=\arg\max_k f_{\theta,k}(z^r)\) và
\(g_{aj}(z)=f_{\theta,a}(z)-f_{\theta,j}(z)\).
\begin{proposition}[Điều kiện đủ của decision preservation]\label{prop}
Giả sử với mỗi \(j\ne a\), \(\ell_{aj}\) là lower bound hợp lệ:
\[
 \forall z\in\Zset_e(z^r):\quad g_{aj}(z)\ge\ell_{aj}.
\]
Nếu \(\min_{j\ne a}\ell_{aj}>0\), mọi điểm trong miền có unique winner \(a\).
Với query \(z_t\in\Zset_e(z^r)\), có thể trả lại lớp \(a\).
\end{proposition}
\begin{proof}
Với mọi đối thủ \(j\), \(f_{\theta,a}(z)-f_{\theta,j}(z)\ge\ell_{aj}>0\).
Vì vậy logit \(a\) lớn hơn từng logit khác trên toàn miền.
\end{proof}
Điều kiện strict-positive là điều kiện đủ; không khẳng định tương đương với
mọi argmax policy có tie. Mệnh đề là cơ sở logic, không phải một novelty theorem.

\subsection{Affine bound cho diagonal region}
Nếu một affine relaxation hợp lệ trên \emph{toàn miền} cho
\(g_{aj}(z)\ge\alpha_{aj}^{\top}z+\beta_{aj}\), thì
\begin{equation}\label{eq:affine}
 \ell_{aj}
 =\alpha_{aj}^{\top}z^r+\beta_{aj}
  -\norm{B_e^\top\alpha_{aj}}_1.
\end{equation}
Công thức dùng support function của \(\ell_\infty\) box.
Với global box radius \(\rho\), đặt \(B_e=\rho I\); penalty là
\(\rho\norm{\alpha_{aj}}_1\). Đặt \(B_e=I\) chỉ cho unit box.
Không lấy affine coefficients từ một miền nhỏ rồi dùng cho miền lớn hơn
mà chưa chứng minh relaxation vẫn hợp lệ.

\subsection{PCA audit và biến thể correlated}
Audit hiện tại dùng raw \(\delta=x_{t+1}-x_t\), train-fit center \(\mu\),
orthonormal PCA columns \(U\) và projection constraints:
\begin{align}
 c&=U^\top(\delta-\mu),&
 v&=(I-UU^\top)(\delta-\mu),\\
 |c_i|&\le b_i,& \norm{v}_2&\le\varepsilon.
\end{align}
Định nghĩa miền bằng chính các projection constraints cho membership xác
định như script. Nó tương đương
\(x=x^r+\mu+Uc+v\), \(U^\top v=0\), với các giới hạn trên.
Nếu affine lower bound trên đúng miền là
\(\alpha^\top x+\beta\), support-function bound là
\begin{equation}
 \ell=\alpha^\top(x^r+\mu)+\beta
 -\norm{\operatorname{diag}(b)U^\top\alpha}_1
 -\varepsilon\norm{(I-UU^\top)\alpha}_2.
\end{equation}
Công thức này là phân tích real arithmetic, chưa phải verifier đã chạy.

Bản draft trước dùng residual coordinate-box \(|v_i|\le\epsilon_i\) và
không có centering. Đó là một biến thể khác, có penalty
\(\sum_i\epsilon_i|\alpha_i|\) khi residual box tự do; không dùng coverage
PCA hiện tại làm evidence cho biến thể đó. Empirical quantile có thể chọn
candidate region; soundness phải đến từ valid bounds và membership,
không phải bản thân quantile hoặc coverage tốt.

\section{Runtime đề xuất và certified reuse gate}
Certificate record cần chứa ít nhất anchor input/class, checkpoint hash,
normalization identity, domain type/parameters, lower margins, verifier
version và numerical assumptions. Không coi semantic event label là proof.
Một selector chọn sai tên event vẫn không gây chấp nhận sai nếu certificate
và membership thực sự đúng; nó có thể làm giảm hiệu quả.

\begin{enumerate}
\item Tạo query đúng preprocessing và tìm certificate tương ứng.
\item Kiểm identity và tính hợp lệ của certificate.
\item Kiểm membership trong đúng miền và mọi stored lower margin dương.
\item Chỉ khi đủ điều kiện mới reuse lớp anchor.
\item Nếu check không đạt: full inference, trả prediction và có thể đặt
trusted anchor mới rồi tạo certificate theo policy.
\end{enumerate}
Không mint trusted anchor chỉ từ một window được reuse. Reset state ở
recording boundary; thay weights, preprocessing hoặc domain identity phải
invalidate certificate liên quan. Construction có thể tốn chi phí rồi được
amortize; thiết kế không mặc định chạy CROWN mỗi window.

Real-arithmetic proof không tự bảo đảm floating-point implementation sound.
G3 phải nêu numerical contract, tie/NaN/infinity handling và cơ sở lower
bounds hợp lệ. Zero observed mismatch chỉ là implementation-test evidence,
không thay soundness argument. Một tolerance tùy ý không phải numerical proof;
nếu nghĩa vụ này chưa giải quyết, guarantee phải ghi điều kiện và G3 vẫn pending.

\section{Dữ liệu và phương pháp đánh giá đã có}
Source dựng 128-sample windows trong từng recording, không nối qua recording.
Split subject-disjoint được lưu như Bảng~\ref{tab:split}; validation dùng
một nhóm ID cố định, chưa có cross-validation.

\begin{table}[ht]
\centering\small
\begin{tabular}{lrrrr}
\toprule
Phạm vi & Train & Val & Test & Tổng\\
\midrule
Mọi stream windows & 9.978 & 2.422 & 5.053 & 17.453\\
Pure K=6 cho G2 & 5.693 & 1.450 & 2.874 & 10.017\\
Subjects & 17 & 4 & 9 & 30\\
\bottomrule
\end{tabular}
\caption{Denominators từ split/training artifacts.}\label{tab:split}
\end{table}

Toàn corpus có 3.587 unlabeled và 1.588 ambiguous windows. Accuracy/F1 hiện
chỉ dùng test pure K=6; không đại diện cho đủ 12 lớp hoặc full stream.
Full test stream 5.053 windows có thể dùng để so decision agreement với
reference, kể cả khi không có ground truth. Train/val/test không được gộp
rồi gọi là test evaluation.

Training dùng Adam, batch 128, tối đa 30 epochs, seeds 42/43/44 và validation
macro-F1 để chọn epoch/seed. Script09 tính lại val/test từ checkpoint.
Không khẳng định test chỉ được truy cập một lần toàn workflow.
Test metrics và exploratory coverage đã được xem trong G1/G2. Các experiment
tiếp theo phải công khai việc đó; không gọi tập test là một confirmatory
holdout hoàn toàn chưa từng được quan sát.

\section{Kết quả baseline được lưu trong project}
\subsection{Backbone classification}
\begin{table}[ht]
\centering\small
\begin{tabular}{rrrrr}
\toprule
Seed & Epoch chọn & Val macro-F1 & Test macro-F1 & Test accuracy\\
\midrule
42 & 4 & 0,9454 & 0,9105 & 0,9085\\
43 & 3 & 0,9418 & 0,8906 & 0,8897\\
44 & 25 & 0,9504 & 0,9354 & 0,9349\\
\bottomrule
\end{tabular}
\caption{Artifacts của các run; test denominator là 2.874 pure K=6 windows.}
\end{table}
Checkpoint seed44 được chọn theo validation. Ba seeds trên cùng split cho
test macro-F1 \(0{,}9122\pm0{,}0183\), accuracy \(0{,}9110\pm0{,}0185\).
Đây là mean và population standard deviation (\code{ddof=0}), không phải
confidence interval. Hai lớp yếu nhất của checkpoint chọn là SITTING
(F1 0,8577) và STANDING (0,8779).

Margin artifact dùng \(f_y-\max_{j\ne y}f_j\) theo ground truth tại từng điểm.
Nó không phải certified radius hoặc lower bound trên một neighborhood.

\subsection{Exploratory region coverage}
Audit có 17.392 adjacent pairs; test gồm 5.034 pairs.
\begin{table}[ht]
\centering\small
\begin{tabularx}{\linewidth}{Yrrr}
\toprule
Candidate & Mọi test pair & Same-activity & Label-change\\
\midrule
Diagonal q75 & 25,53\% & 45,81\% & 0,21\%\\
Diagonal q90 & 31,94\% & 52,54\% & 2,26\%\\
PCA, 108 components & 35,42\% & -- & --\\
\bottomrule
\end{tabularx}
\caption{Membership coverage, chưa phải certified reuse. Dấu --: summary PCA
không lưu riêng metric đó.}
\end{table}
PCA hơn diagonal q90 3,48 điểm phần trăm; comparison chưa matched budget/
coverage. Same-activity là hai pure windows cùng majority label; nhóm
label-change dùng majority labels khác nhau. Adjacent-pair coverage không
cho reuse rate với anchor giữ qua nhiều bước.
Các số audit là bản sau index fix; không dùng pre-fix values.

\subsection{Forward-only CPU baseline}
\begin{table}[ht]
\centering\small
\begin{tabular}{lrr}
\toprule
Metric & Batch 1 (ms/window) & Batch 32 (ms/batch)\\
\midrule
Mean & 0,2511 & 0,9663\\
p50 & 0,2357 & 0,9323\\
p95 & 0,3313 & 1,1700\\
\bottomrule
\end{tabular}
\caption{Saved latency artifact, CPU, sáu Torch threads.}
\end{table}
PyTorch 2.14.1+cpu, \code{model.eval()}, \code{no_grad()}, tensor đã normalize,
resident RAM. Batch 1: 200 warmups và 1.000 timing runs. Timing chỉ
\code{model(x)}, không bao data loading, windowing, normalization, event,
membership, certificate/cache/rebuild. CPU model và OS của run chưa được
lưu. Chưa đo lại trong lượt biên soạn; không chuyển các số này sang thiết bị
edge hoặc latency end-to-end.

\section{Cost accounting và điều kiện break-even}
Một mô hình trung bình đơn giản, khi forward cost tương đối đồng nhất, là
\begin{align}
 \mathbb{E}[T]&\approx T_e+T_m+(1-r)T_f+qT_b,\\
 \text{saving}&:\quad rT_f>T_e+T_m+qT_b.
\end{align}
\(T_e\): mean event cost; \(T_m\): mean lookup/membership/validity cost;
\(r\): reuse fraction trên denominator khai báo; \(T_f\): mean forward;
\(q\): rebuild count/window; \(T_b\): mean cost trên các rebuild.
Initial build phải nằm trong count hoặc được tính thêm \(T_{\rm init}/N\)
trên horizon \(N\). Preprocessing chung vẫn được timing cho cả hai pipeline.

Với mean baseline \(T_f=0{,}2511\) ms và \emph{giả sử} \(r=0{,}5\), ngân sách
overhead còn \(0{,}12555\) ms, tức 125,55\(\,\mu\)s/window. Đây là sensitivity
calculation, không phải reuse rate đo được. Không thay mean bằng p50 trong
công thức kỳ vọng; p95 phải lấy từ runtime trace.

Một stream có khoảng cách window 1.280 ms; forward mean bằng khoảng 0,0196\%
khoảng này. Baseline hiện không thể hiện áp lực deadline realtime single-stream
trên CPU đó. Giá trị về compute/đa stream hoặc thiết bị chậm hơn cần được
đo riêng. Net saving cuối cùng dựa trên end-to-end traces, không chỉ mô hình
trung bình này.

\section{Kế hoạch G3--G6 và tiêu chí hoàn thành}
\begin{longtable}{p{0.08\linewidth}p{0.31\linewidth}p{0.52\linewidth}}
\toprule
Gate & Đầu ra & Điều kiện kiểm chứng\\
\midrule
\endhead
G3.1 & Domain contract & Raw/normalized mapping, preprocessing, diagonal
region, model/anchor/domain identity nhất quán.\\
G3.2 & Minimal verifier & Lower bounds hợp lệ trên toàn diagonal region của
frozen CNN; numerical assumptions rõ.\\
G3.3 & Certified reuse gate & Membership và identity checks; fail fallback;
anchor lifecycle; trace và implementation tests.\\
G3.4 & Component instrumentation & Đo construction, membership, fallback,
inference và initial/rebuild costs. Chưa tự kết luận net saving.\\
G4 & Structured ablation & Giữ model/bound method khi đổi region; giữ region
khi đổi bound method; so sánh công bằng.\\
G5 & End-to-end evaluation & Mean/p50/p95, reuse, mismatches, rebuild,
memory và full-pipeline baseline cùng môi trường.\\
G6 & Báo cáo và reproducibility & Claims, evidence, units, denominators,
limitations và scripts có thể tái kiểm.\\
\bottomrule
\end{longtable}
Baseline tối thiểu theo reading plan: full inference, confidence gate,
global Lipschitz, diagonal CROWN và structured CROWN, tùy feasibility.
Tune policy trên validation; test-stream evaluation dùng policy đã cố định.
Không dùng ground-truth event ở runtime mà bỏ qua cost/availability. Nếu không
có query nào được reuse, empirical agreement và mismatch rate trên reused
subset là N/A, không báo measured agreement 100\%.

\section{Giới hạn và threats to validity}
\textbf{Internal validity.} Chưa tái lập saved runs trong lượt viết; chưa
có verifier, gate hoặc numerical soundness evidence. Bản PCA trước đây có
domain mismatch; contract mới sửa mô tả, không có nghĩa code verifier đã tồn tại.

\textbf{Construct validity.} Coverage khác certified reuse; ground-truth
accuracy khác reference agreement; verification speed khác inference speed.
Không suy certificate từ high confidence hoặc observed zero mismatch.

\textbf{External validity.} Một dataset, một fixed subject split, một K=6 CNN
và một saved CPU setting chưa hỗ trợ generalization sang 12 classes,
dataset khác, device khác hoặc production. Biến thiên seed không phải
subject confidence interval.

\textbf{Systems validity.} Forward baseline dùng tensor resident lặp lại;
event projection, cache, rebuild, preprocessing và memory có thể triệt tiêu
saving. Chi phí certificate ban đầu và horizon amortization phải được khai báo.

\textbf{Literature validity.} Chưa hoàn tất P01--P23 và chưa chứng minh novelty.
LipCache dùng local PDF với provenance đã ghi; không có fresh successful HTML
fetch trong lượt 22/10. STBP local hash không chứng minh đồng nhất byte với
remote PDF hiện tại.

\section{Kết luận và khả năng tái kiểm}
ECDR có bài toán và điều kiện đủ rõ để thử nghiệm khả thi, cùng data/model
baseline đã lưu. Phần nghiên cứu trọng yếu vẫn mở: verifier sound trên đúng
domain, giá trị event structure và net inference saving.
Trạng thái theo gate record: G0 PASS, G1 COMPLETE\_WITH\_REVISIONS, G2 PASS;
G3--G6 pending. Các trạng thái này không phải phần trăm research claims đã
được xác nhận. Chưa có claim production hoặc method acceleration đã chứng minh.

Nguồn số liệu: \code{data/g2_metrics.json}, \code{data/g2_training.json},
\code{data/g2_model_summary.json}, \code{data/g2_cpu_latency.json},
\code{data/event_audit.json}, \code{data/split_audit.json}.
Source chính: \code{scripts/backbone.py}, scripts03--10 và configs YAML.
Tài liệu đồng bộ: \code{PROJECT_BRIEF.md}, \code{04_THEORY.md},
\code{07_G3_DOMAIN_CONTRACT.md}, \code{DOCS_SYNC_2026-10-22.md}.

Checkpoint \code{backbone_best.pt} có SHA-256:
\begin{quote}\small\ttfamily
17e7cfa3877134feaa1392d989ffb824a68de266338da69bc8e43c0ae59bcd1d
\end{quote}
Hash nhận diện file; chưa load weights để tái tạo metrics trong lượt này.
Báo cáo là tài liệu làm việc có phân tách evidence/proposal, chưa có author
roster, venue, deadline hoặc approval record để tự bổ sung.

\begin{thebibliography}{10}
\bibitem{hapt}
J. Reyes-Ortiz, D. Anguita, L. Oneto, X. Parra.
\emph{Smartphone-Based Recognition of Human Activities and Postural Transitions}.
UCI Machine Learning Repository, 2015. DOI:\href{https://doi.org/10.24432/C54G7M}{10.24432/C54G7M}.

\bibitem{crown}
H. Zhang, T.-W. Weng, P.-Y. Chen, C.-J. Hsieh, L. Daniel.
\emph{Efficient Neural Network Robustness Certification with General Activation Functions}.
NeurIPS 31, 2018.
\url{https://proceedings.neurips.cc/paper/2018/hash/d04863f100d59b3eb688a11f95b0ae60-Abstract.html}.

\bibitem{lirpa}
K. Xu, Z. Shi, H. Zhang, Y. Wang, K.-W. Chang, M. Huang,
B. Kailkhura, X. Lin, C.-J. Hsieh.
\emph{Automatic Perturbation Analysis for Scalable Certified Robustness and Beyond}.
NeurIPS 33, 2020.
\url{https://papers.nips.cc/paper/2020/hash/0cbc5671ae26f67871cb914d81ef8fc1-Abstract.html}.

\bibitem{api}
auto-LiRPA contributors. \emph{API Usage}.
Tài liệu kỹ thuật, truy cập 22/10/2026.
\url{https://auto-lirpa.readthedocs.io/en/latest/api.html}.

\bibitem{wei}
T. Wei, C. Liu.
\emph{Online Verification of Deep Neural Networks under Domain Shift or Network Updates}.
arXiv:2106.12732,v1(2021),v2(2023).
\url{https://arxiv.org/abs/2106.12732}.

\bibitem{noscope}
D. Kang, J. Emmons, F. Abuzaid, P. Bailis, M. Zaharia.
\emph{NoScope: Optimizing Neural Network Queries over Video at Scale}.
PVLDB 10(11):1586--1597,2017.
DOI:\href{https://doi.org/10.14778/3137628.3137664}{10.14778/3137628.3137664}.

\bibitem{deltacnn}
M. Parger, C. Tang, C. D. Twigg, C. Keskin, R. Wang, M. Steinberger.
\emph{DeltaCNN: End-to-End CNN Inference of Sparse Frame Differences in Videos}.
CVPR, 2022,12497--12506.
\url{https://openaccess.thecvf.com/content/CVPR2022/html/Parger_DeltaCNN_End-to-End_CNN_Inference_of_Sparse_Frame_Differences_in_Videos_CVPR_2022_paper.html}.

\bibitem{lipcache}
Z. Xiang, Y. Chen, F. Ying, B. Zhou, H. Zhao, S. Dustdar.
\emph{LipCache: A Local Inference Proxy with Certified Caching for Edge Image Classification Service}.
arXiv:2608.13144v1,2026.
\url{https://arxiv.org/abs/2608.13144}.

\bibitem{stbp}
S. Varghese, M. Wicker, A. Lomuscio.
\emph{Hybrid Robustness Verification for Spatio-Temporal Neural Networks}.
arXiv:2606.09746v1,2026.
\url{https://arxiv.org/abs/2606.09746}.

\bibitem{franco}
N. Franco, J. Spiegelberg, J. M. Lorenz, S. G\"unnemann.
\emph{Guaranteeing Robustness Against Real-World Perturbations In Time Series
Classification Using Conformalized Randomized Smoothing}.
UAI,PMLR 244:1371--1388,2024.
\url{https://proceedings.mlr.press/v244/franco24a.html}.
\end{thebibliography}
\end{document}
